\documentclass[11pt]{article}
\usepackage[margin=0.85in]{geometry}
\usepackage[T1]{fontenc}
\usepackage{amsmath,amssymb,booktabs,graphicx,xurl,hyperref}
\graphicspath{{outputs/task1/}{outputs/task2/}{outputs/task3/}{outputs/task4/}}

\title{CS 498 HW3 Report\\Center-Based 3D Detection from LiDAR}
\author{Your Name (NetID)}
\date{Fall 2026}

\begin{document}
\maketitle

\section{BEV representation}
State the LiDAR coordinate convention and explain the four BEV channels.
Include \path{bev_channels.png} and \path{bev_boxes.png}. Report tensor shape,
metric resolution, occupancy, and channel ranges. Explain how you confirmed
that the vehicle boxes align with occupied LiDAR structure.

\section{Center-based targets}
Explain the center heatmap, sub-cell offset, height, log-dimension, and
sine/cosine yaw targets. Include \path{target_roundtrip.png}. Report the number
of encoded boxes, heatmap maximum, and metric center error after decoding.

\section{Detector and optimization}
Describe the supplied residual backbone and the five output heads, including
their tensor shapes. Give the focal and masked regression loss equations.
Include \path{training_curves.png}; state final hyperparameters, parameter
count, best validation epoch, loss behavior, and validation AP@0.5. Discuss
which loss term was most difficult to optimize and why.

\section{Inference and evaluation}
Explain local-maximum suppression, top-k selection, metric box decoding,
rotated NMS, and validation threshold selection. Explain why AP uses the full
ranked prediction list while precision and recall use the selected operating
threshold. Include the validation PR curve and at least two fixed validation
visualizations with ground truth and predictions. Report validation precision,
recall, and AP@0.5, then report the
single final test AP@0.5. Discuss one false positive, one false negative, and
what the PR-curve shape reveals.

\section{Limitations and next step}
Discuss at least two limitations involving BEV resolution, single-sweep input,
class coverage, data imbalance, or range. Propose one concrete next experiment.

\section*{Optional bonus}
If you attempt tracking, state the Kalman state and motion/measurement models,
association cost and gate, and birth/death settings. Name your submitted
\path{tracking_comparison.mp4}; every predicted ID must keep one stable color.
Report results at rotated BEV IoU 0.5:

\begin{center}
\begin{tabular}{@{}rrrr@{}}
\toprule
MOTA & MOTP (mean IoU) & IDF1 & ID switches \\
\midrule
-- & -- & -- & -- \\
\bottomrule
\end{tabular}
\end{center}

Explain one failure case involving association, occlusion, track management,
or detector errors.

If you attempt multi-class detection, identify the optional public companion
labels and explain how their rows align to the unchanged core BEV tensors.
Report Car, Pedestrian, and Cyclist validation AP at rotated BEV IoU 0.5,
show a per-class PR curve, and analyze at least two fixed qualitative
validation views. State how you selected the multiclass checkpoint and any
per-class operating thresholds on validation. If you also report final test
metrics, confirm one evaluation after choices were frozen and discuss the
small Cyclist test count. Include all bonus code with the submission.

\section*{AI-use declaration}
Check one statement.
\begin{itemize}
  \item[$\square$] I did not use generative AI for this homework.
  \item[$\square$] I did not use generative AI for any required part. I used it
  only for an optional bonus marked ``Gen-AI leverage allowed.'' Tool/model,
  purpose, and affected files: \underline{\hspace{1.4in}}
\end{itemize}
I confirm that generative AI did not write or complete my required code,
derive my required answers, generate my analysis, or write my required report.

\end{document}
